\begin{abstract}
As multi-agent societies powered by Large Language Models (LLMs) are increasingly deployed for social simulation, decision support, and collective deliberation, a critical question emerges: do autonomous agents seek empirical truth, or do they prioritize social alignment and conversational comfort? In human communication, beliefs often diffuse based on interpersonal rapport and authority rather than objective verification. In this work, we introduce an empirical framework to quantify the \textit{Truth Discount}---the systematic penalty that empirical veracity incurs relative to socially endorsed claims within conversational agent networks. We present a strictly counterbalanced, paired-world simulation paradigm in which matched agent communities (identical in topology, demographics, and initial conditions) are exposed to either a verifiable truth or a semantically matched rumor. Agents navigate sequential multi-agent interactions while deciding whether to expend private budget on costly empirical verification. By evaluating source authority, social closeness, and baseline reliability, we isolate the conditions under which agents voluntarily verify ambiguous claims versus deferring to peer testimony. Our framework establishes standardized behavioral metrics for belief endorsement, transmission fidelity, and voluntary verification in LLM societies.
\end{abstract}
